\SourceLeaf{47}\emph{p.} 109. Quia vero evolutio potestatis binomii a theoria numerorum satis aliena esse videbatur, aliam demonstrationem ill. Euler investigavit quae exstat \emph{Comment. nov. Petr. T.} VII \emph{p.} 70, atque cum ea quam nos art. praec. exposuimus prorsus convenit. In sequentibus adhuc aliae quaedam se nobis offerent. Hoc loco unam superaddere liceat, quae similibus principiis innititur, uti prima ill. Euleri. Propositio sequens, cuius casus tantum particularis est theorema nostrum, etiam ad alias investigationes infra adhibebitur.

\Article{51}

\emph{Polynomii $a+b+c+\mathit{etc.}$ potestas $p^{ta}$ secundum modulum $p$ est}

\[\equiv a^p+b^p+c^p+\mathit{etc.}\]

\noindent\emph{siquidem $p$ est numerus primus.}

\emph{Demonstr.} Constat potestatem $p^{\mathrm{tam}}$ polynomii $a+b+c+\text{etc.}$ esse compositam e partibus formae $\varkappa a^\alpha b^\beta c^\gamma$ etc. ubi $\alpha+\beta+\gamma\text{ etc.}=p$, et $\varkappa$ designat, quot modis $p$ res, quarum $\alpha,\beta,\gamma$ etc. respective sunt $=a,b,c$ etc. permutari possint. At supra art. 41 ostendimus, hunc numerum semper esse per $p$ divisibilem, nisi omnes res sint aequales, \emph{i. e.} nisi aliquis numerorum $\alpha,\beta,\gamma$ etc. sit $=p$ reliqui vero $=0$. Unde sequitur omnes ipsius $(a+b+c+\text{etc.})^p$ partes, praeter has $a^p,b^p,c^p$ etc., per $p$ divisibiles esse; quae igitur quando de congruentia secundum modulum $p$ agitur, tuto omitti poterunt, fietque

\[(a+b+c+\mathrm{etc.})^p\equiv a^p+b^p+c^p+\mathrm{etc.}\quad\text{\emph{Q. E. D.}}\]

Quodsi iam omnes quantitates $a,b,c$ etc. $=1$ ponuntur, numerusque earum $=k$, fiet $k^p\equiv k$ uti in art. praec.

\begin{center}\emph{Quot numeris respondeant periodi, in quibus terminorum multitudo est divisor datus numeri $p-1$.}\end{center}

\Article{52}

Quoniam igitur alii numeri quam qui sunt divisores ipsius $p-1$ nequeunt esse exponentes potestatum infimarum ad quas evecti numeri aliqui unitati congrui fiunt, quaestio sese offert, num omnes ipsius $p-1$ divisores ad hoc sint idonei, atque, quando omnes numeri per $p$ non divisibiles secundum exponentem infimae suae potestatis unitati congruae classificentur, quot ad singulos exponentes sint perventuri. Ubi statim observare convenit, sufficere, si omnes numeri
\SourceLeaf{48}positivi ab $1$ usque ad $p-1$ considerentur; manifestum enim est, numeros congruos ad eandem potestatem elevari debere, quo unitati fiant congruae, adeoque numerum quemcunque ad eundem exponentem esse referendum ad quem residuum suum minimum positivum. Quocirca in id nobis erit incumbendum, ut quomodo hoc respectu numeri $1,2,3\GaussFourDots p-1$ inter singulos factores numeri $p-1$ distribuendi sint eruamus. Brevitatis gratia, si $d$ est unus e divisoribus numeri $p-1$ (ad quos etiam $1$ et $p-1$ referendi), per $\psi d$ designabimus multitudinem numerorum positivorum ipso $p$ minorum quorum potestas $d^{\mathrm{ta}}$ est infima unitati congrua.

\Article{53}

Quo facilius haec disquisitio intelligi possit, exemplum apponimus. Pro $p=19$ distribuentur numeri $1,2,3\GaussFourDots18$ inter divisores numeri $18$ hoc modo:

\[
\begin{array}{r|l}
1 & \phantom{1}1.\\
2 & 18.\\
3 & \phantom{1}7,\ 11.\\
6 & \phantom{1}8,\ 12.\\
9 & \phantom{1}4,\ 5.\ 6,\ 9,\ 16,\ 17.\\
18 & \phantom{1}2,\ 3,\ 10,\ 13.\ 14,\ 15.
\end{array}
\]

In hoc igitur casu fit $\psi1=1$, $\psi2=1$, $\psi3=2$, $\psi6=2$, $\psi9=6$, $\psi18=6$. Ubi exigua attentio docet, totidem ad quemvis exponentem pertinere, quot dentur numeri hoc non maiores ad ipsumque primi, sive esse in hoc certe casu, retento signo art. 39, $\psi d=\phi d$. Hanc autem observationem generaliter veram esse ita demonstramus.

I. Si numerus aliquis habetur, $a$, ad exponentem $d$ \emph{pertinens} (\emph{i. e.} cuius potestas $d^{\mathrm{ta}}$ unitati congrua, omnes inferiores incongruae), omnes huius potestates $aa,a^3,a^4\GaussFourDots a^d$ sive ipsarum residua minima proprietatem priorem etiam possidebunt (ut potestas ipsarum $d^{\mathrm{ta}}$ unitati sit congrua) et quum hoc ita etiam exprimi possit, residua minima numerorum $a,aa,a^3\GaussFourDots a^d$ (quae omnia sunt diversa) esse radices congruentiae $x^d\equiv1$, haec autem plures quam $d$ radices diversas habere nequeat, manifestum est, praeter numerorum $a,aa,a^3\GaussFourDots a^d$ residua minima alios numeros inter $1$ et $p-1$ incl. non dari quorum potestates ex\SourceLeaf{49}ponentis $d$ congruae sint unitati. Hinc patet omnes numeros ad exponentem $d$ pertinentes inter residua minima numerorum $a,aa,a^3\GaussFourDots a^d$ reperiri. Quales vero sint, quantaque eorum multitudo, ita definitur. Si $k$ est numerus ad $d$ primus, omnes potestates ipsius $a^k$, quarum exponentes $<d$, unitati non erunt congrui: esto enim $\frac{1}{k}\,(\operatorname{mod.}d)\equiv m$ (vid. art. 31) eritque $a^{km}\equiv a$; quare si potestas $e^{\mathrm{ta}}$ ipsius $a^k$ unitati esset congrua atque $e<d$, foret etiam $a^{kme}\equiv1$ et hinc $a^e\equiv1$ contra hyp. Hinc manifestum est, residuum minimum ipsius $a^k$ ad exponentem $d$ pertinere. Si vero $k$ divisorem aliquem, $\delta$, cum $d$ communem habet, ipsius $a^k$ residuum minimum ad exponentem $d$ non pertinet; quoniam tum potestas ${\frac{d}{\delta}}^{\mathrm{ta}}$ iam unitati fit congrua (erit enim $\frac{kd}{\delta}$ per $d$ divisibilis, sive $\equiv0\,(\operatorname{mod.}d)$ adeoque $a^{\frac{kd}{\delta}}\equiv1$). Hinc colligitur, totidem numeros ad exponentem $d$ pertinere quot numerorum $1,2,3\GaussFourDots d$ ad $d$ sint primi. At memorem esse oportet, hanc conclusionem innixam esse suppositioni, unum numerum $a$ iam haberi ad exponentem $d$ pertinentem. Quamobrem dubium remanet, fierine possit ut ad aliquem exponentem nullus omnino numerus pertineat; conclusioque eo limitatur ut $\psi d$ sit vel $=0$ vel $=\phi d$.

\Article{54}

II. Iam sint omnes divisores numeri $p-1$ hi: $d,d',d''$ etc. eritque, quia omnes numeri $1,2,3\GaussFourDots p-1$ inter hos sunt distributi,

\[\psi d+\psi d'+\psi d''+\text{etc.}=p-1\]

At in art. 40 demonstravimus esse

\[\phi d+\phi d'+\phi d''+\text{etc.}=p-1\]

atque ex art. praec. sequitur $\psi d$ ipsi $\phi d$ aut aequalem aut ipso minorem esse, maiorem esse non posse, similiterque de $\psi d'$ et $\phi d'$, etc. Si itaque aliquis terminus ex his $\psi d,\psi d',\psi d''$ etc. termino respondente ex his $\phi d,\phi d',\phi d''$, esset minor (sive etiam plures) illorum summa summae horum aequalis esse non posset. Unde tandem concludimus $\psi d$ \emph{ipsi $\phi d$ semper esse aequalem}, adeoque a magnitudine ipsius $p-1$ non pendere.

\Article{55}

Maximam autem attentionem meretur casus particularis propositionis prae\SourceLeaf{50}cedentis scilicet \emph{semper dari numeros quorum nulla potestas inferior quam $p-1^{ta}$ unitati congrua}, et quidem totidem inter $1$ et $p-1$ quot infra $p-1$ sint numeri ad $p-1$ primi. Cuius theorematis demonstratio quum minime tam obvia sit quam primo aspectu videri possit, propter theorematis dignitatem liceat aliam adhuc adiicere a praecedente aliquantum diversam, quandoquidem methodorum diversitas ad res obscuriores illustrandas plurimum conferre solet. Resolvatur $p-1$ in factores suos primos fiatque $p-1=a^\alpha b^\beta c^\gamma$ etc., designantibus $a,b,c$ etc. numeros primos inaequales. Tum theorematis demonstrationem per sequentia absolvemus:

I. Semper inveniri posse numerum $A$ (aut plures) ad exponentem $a^\alpha$ pertinentem, similiterque numeros $B,C$ etc. ad exponentes $b^\beta,c^\gamma$ etc. respective pertinentes.

II. Productum ex omnibus numeris $A,B,C$ etc. (sive huius producti residuum minimum) ad exponentem $p-1$ pertinere. Haec autem ita demonstramus.

I. Sit $g$ numerus aliquis ex his $1,2,3\ldots p-1$, congruentiae $x^{\frac{p-1}{a}}\equiv1\,(\operatorname{mod.}p)$ \emph{non} satisfaciens, omnes enim hi numeri congruentiae huic, cuius gradus $<p-1$, satisfacere nequeunt. Tum dico si potestas ${\frac{p-1}{a^\alpha}}^{\mathrm{ta}}$ ipsius $g$ ponatur $\equiv h$, hunc numerum, sive eius residuum minimum ad exponentem $a^\alpha$ pertinere.

Namque patet potestatem $a^{\alpha\,\mathrm{tam}}$ ipsius $h$ congruam fore potestati $p-1^{\mathrm{tae}}$ ipsius $g$ \emph{i. e.} unitati, potestas vero $a^{\alpha-1\,\mathrm{ta}}$ ipsius $h$ congrua erit potestati ${\frac{p-1}{a}}^{\mathrm{tae}}$ ipsius $g$, \emph{i. e.} unitati erit incongrua, multoque minus potestates $a^{\alpha-2},\ a^{\alpha-3\,\mathrm{tae}}$ etc. ipsius $h$ unitati congruae esse possunt. At exponens infimae potestatis ipsius $h$, unitati congruae, sive exponens ad quam pertinet $h$, numerum $a^\alpha$ metiri debet (art. 48). Quare quum $a^\alpha$ per alios numeros divisibilis non sit quam per se ipsum, atque per inferiores ipsius $a$ potestates, necessario $a^\alpha$ erit exponens ad quem $h$ pertinet. \mbox{\emph{Q. E. D.}} Per similem methodum demonstratur, dari numeros ad exponentes $b^\beta,c^\gamma$ etc. pertinentes.

II. Si supponimus, productum ex omnibus $A,B,C$ etc. non ad exponentem $p-1$, sed ad minorem $t$ pertinere, $t$ ipsum $p-1$ metietur (art. 48), sive erit $\frac{p-1}{t}$ integer unitate maior. Facile autem perspicitur, hunc quotientem vel esse unum e numeris primis $a,b,c$ etc. vel saltem per aliquem eorum divisibilem (art. 17), \emph{ex. gr.} per $a$, de reliquis enim simile est ratiocinium.
